% =========================================================================
\section{Operadores diferenciales: divergencia, rotacional y laplaciano}\label{sec:operadores}
% =========================================================================

% -------------------------------------------------------------------------
\subsection{El operador nabla}\label{sec:nabla}
% -------------------------------------------------------------------------
Todo el cálculo diferencial vectorial se organiza alrededor del operador
simbólico
\begin{equation}\label{eq:nabla}
  \nab=\pd{}{x}\,\ux+\pd{}{y}\,\uy+\pd{}{z}\,\uz=\sum_{i=1}^{3}\uh{i}\,\partial_i ,
\end{equation}
que se combina con los campos de las tres maneras en que un vector puede
combinarse con otros objetos (\cref{obs:nabla}):
\begin{center}
\renewcommand{\arraystretch}{1.25}
\begin{tabular}{@{}lcll@{}}
$\vc{A}\,f$ & $\longrightarrow$ & $\nab f$ & (gradiente), aplicado a un campo escalar;\\
$\vc{A}\cdot\vc{B}$ & $\longrightarrow$ & $\nab\cdot\vc{B}$ & (divergencia), aplicado a un campo vectorial;\\
$\vc{A}\times\vc{B}$ & $\longrightarrow$ & $\nab\times\vc{B}$ & (rotacional), aplicado a un campo vectorial.
\end{tabular}
\end{center}
El gradiente ya se estudió en la \cref{sec:campos}; veamos los otros dos.

% -------------------------------------------------------------------------
\subsection{Divergencia}\label{sec:divergencia}
% -------------------------------------------------------------------------
\begin{definicion}[Divergencia]\label{def:divergencia}
Sea $\vc{F}=F_x\,\ux+F_y\,\uy+F_z\,\uz$ un campo vectorial derivable. Su
\emph{divergencia} es el campo escalar
\begin{equation}\label{eq:div}
  \nab\cdot\vc{F}(\vc{r})=\pd{F_x}{x}+\pd{F_y}{y}+\pd{F_z}{z}
  =\nab\cdot\bigl(F_x\,\ux+F_y\,\uy+F_z\,\uz\bigr).
\end{equation}
\end{definicion}

\emph{¿Qué mide?} La divergencia dice cuánto «sale» o «entra» del campo en un
entorno infinitesimal del punto. Si pensamos en $\vc{F}$ como el campo de
velocidades de un fluido, $\nab\cdot\vc{F}(\vc{r})$ mide la tendencia del fluido
a expandirse o a contraerse en $\vc{r}$. En la \cref{sec:gauss} daremos la
definición intrínseca (independiente de las coordenadas) como densidad de flujo
por unidad de volumen:
\begin{center}
\begin{tabular}{@{}lll@{}}
$\nab\cdot\vc{F}(\vc{r})>0$ & $\Rightarrow$ & el campo \emph{se expande}: $\vc{r}$ es una \emph{fuente};\\
$\nab\cdot\vc{F}(\vc{r})<0$ & $\Rightarrow$ & el campo \emph{se contrae}: $\vc{r}$ es un \emph{sumidero};\\
$\nab\cdot\vc{F}(\vc{r})=0$ & $\Rightarrow$ & campo \emph{solenoidal} (ni se expande ni se contrae).
\end{tabular}
\end{center}

\begin{ejemplo}[Divergencia de campos de referencia]\label{ej:div}
Consideremos los campos de las \cref{fig:radial,fig:campos_div}:
\begin{enumerate}[label=(\roman*)]
  \item $\vc{F}(\vc{r})=x\,\ux+y\,\uy+z\,\uz=\vc{r}$ (campo radial):
        $\nab\cdot\vc{F}=\pd{x}{x}+\pd{y}{y}+\pd{z}{z}=1+1+1=3>0$. Se expande en
        todos los puntos, porque las flechas crecen alejándose del origen.
  \item $\vc{F}_1(\vc{r})=\ux+\uy$ (campo uniforme): $\nab\cdot\vc{F}_1=0$.
  \item $\vc{F}_2(\vc{r})=x\,\ux$: $\nab\cdot\vc{F}_2=\pd{x}{x}=1>0$.
  \item $\vc{F}_3(\vc{r})=y\,\ux$: $\nab\cdot\vc{F}_3=\pd{y}{x}=0$. El módulo
        crece con $y$, pero el campo apunta según $\ux$ y no hay variación del
        módulo en esa dirección.
\end{enumerate}
\end{ejemplo}

\begin{observacion}[Cuándo hay divergencia]\label{obs:cuando-div}
Para que un campo tenga divergencia no nula, la componente $F_i$ tiene que
\emph{variar con respecto a la coordenada $x_i$ correspondiente}: el módulo
del campo debe cambiar en la dirección hacia la que apunta. Si el módulo solo
cambia en una dirección perpendicular al campo (como en $\vc{F}_3=y\,\ux$), la
divergencia es nula: lo que entra por un lado, sale por el otro.
\end{observacion}

% Figura: campo radial F(r)= x ux + y uy + z uz en 3D, en tres planos de
% corte y en el plano xy.
\newcommand{\flechaIII}[6]{%  componentes de posición #1,#2,#3 y del campo #4,#5,#6
  \pgfmathsetmacro{\fm}{sqrt((#4)*(#4)+(#5)*(#5)+(#6)*(#6))}%
  \ifdim\fm pt>0.01pt
    \pgfmathsetmacro{\ft}{\fm/3.465}\mapcolor{\ft}%
    \pgfmathsetmacro{\hl}{2.0+3.4*\ft}%
    \draw[-{Stealth[length=\hl pt,width=\hl pt]}, line width={0.35+0.5*\ft}, mapc]
      (#1-0.17*#4,#2-0.17*#5,#3-0.17*#6) -- ++(0.34*#4,0.34*#5,0.34*#6);
  \fi}
\begin{figure}[ht]
\centering
\tdplotsetmaincoords{68}{128}
\begin{minipage}[t]{0.325\textwidth}\centering
\begin{tikzpicture}[tdplot_main_coords, scale=0.62, baseline=(current bounding box.north)]
  \draw[eje] (0,0,0) -- (3.0,0,0) node[below left] {\small $x$};
  \draw[eje] (0,0,0) -- (0,3.0,0) node[right] {\small $y$};
  \draw[eje] (0,0,0) -- (0,0,3.0) node[above] {\small $z$};
  \foreach \z in {-2,-1,0,1,2}{\foreach \y in {2,1,0,-1,-2}{\foreach \x in {-2,-1,0,1,2}{%
    \flechaIII{\x}{\y}{\z}{\x}{\y}{\z}}}}
\end{tikzpicture}\\[2pt]{\small\scshape (a)}
\end{minipage}\hfill
\begin{minipage}[t]{0.325\textwidth}\centering
\begin{tikzpicture}[tdplot_main_coords, scale=0.62, baseline=(current bounding box.north)]
  \fill[grispizarra!22, opacity=0.55] (-2.7,-2.7,0) -- (2.7,-2.7,0) -- (2.7,2.7,0) -- (-2.7,2.7,0) -- cycle;
  \fill[grispizarra!22, opacity=0.55] (0,-2.7,-2.7) -- (0,2.7,-2.7) -- (0,2.7,2.7) -- (0,-2.7,2.7) -- cycle;
  \fill[grispizarra!22, opacity=0.55] (-2.7,0,-2.7) -- (2.7,0,-2.7) -- (2.7,0,2.7) -- (-2.7,0,2.7) -- cycle;
  \draw[eje] (0,0,0) -- (3.0,0,0) node[below left] {\small $x$};
  \draw[eje] (0,0,0) -- (0,3.0,0) node[right] {\small $y$};
  \draw[eje] (0,0,0) -- (0,0,3.0) node[above] {\small $z$};
  \foreach \x in {-2,-1,0,1,2}{\foreach \y in {-2,-1,0,1,2}{\flechaIII{\x}{\y}{0}{\x}{\y}{0}}}
  \foreach \y in {-2,-1,0,1,2}{\foreach \z in {-2,-1,0,1,2}{\flechaIII{0}{\y}{\z}{0}{\y}{\z}}}
  \foreach \x in {-2,-1,0,1,2}{\foreach \z in {-2,-1,0,1,2}{\flechaIII{\x}{0}{\z}{\x}{0}{\z}}}
\end{tikzpicture}\\[2pt]{\small\scshape (b)}
\end{minipage}\hfill
\begin{minipage}[t]{0.325\textwidth}\centering
\begin{tikzpicture}[x=0.50cm, y=0.50cm, baseline=(current bounding box.north)]
  \marcoD{-2.7}{2.7}{-2.7}{2.7}{-2,-1,0,1,2}{-2,-1,0,1,2}
  \ejesD{-2.7}{2.7}{-2.7}{2.7}
  \clipD{-2.7}{2.7}{-2.7}{2.7}{\flechas{\x}{\y}{-2.4,-1.8,...,2.4}{-2.4,-1.8,...,2.4}{0.5}{3.4}}
\end{tikzpicture}\\[2pt]{\small\scshape (c)}
\end{minipage}
\caption{Campo radial $\vc{F}(\vc{r})=x\,\ux+y\,\uy+z\,\uz=\vc{r}$.
(a)~Representación tridimensional: las flechas se alejan del origen y su
módulo crece linealmente con la distancia (color de índigo a ámbar).
(b)~Cortes por los tres planos coordenados. (c)~Corte por el plano $z=0$.
Es el prototipo de campo con divergencia positiva ($\nab\cdot\vc{F}=3$):
el origen es una fuente.}
\label{fig:radial}
\end{figure}

% Figura: tres campos planos de referencia: uniforme, F=x ux (divergencia
% positiva) y F=y ux (cizalla: rotacional no nulo, divergencia nula).
\begin{figure}[ht]
\centering
\begin{minipage}[t]{0.325\textwidth}\centering
\begin{tikzpicture}[x=0.68cm, y=0.68cm, baseline=(current bounding box.north)]
  \marcoD{-2.7}{2.7}{-2.7}{2.7}{-2,-1,0,1,2}{-2,-1,0,1,2}
  \ejesD{-2.7}{2.7}{-2.7}{2.7}
  \foreach \x in {-2.4,-1.8,...,2.4}{\foreach \y in {-2.4,-1.8,...,2.4}{%
    \draw[-{Stealth[length=4.5pt,width=4.5pt]}, line width=0.7pt, azulprof]
      (\x-0.17,\y-0.17) -- ++(0.34,0.34);}}
\end{tikzpicture}\\[3pt]{\small $\vc{F}_1=\ux+\uy$}
\end{minipage}\hfill
\begin{minipage}[t]{0.325\textwidth}\centering
\begin{tikzpicture}[x=1.12cm, y=1.12cm, baseline=(current bounding box.north)]
  \marcoD{0}{3.2}{0}{3.2}{0,1,2,3}{0,1,2,3}
  \ejesD{0}{3.2}{0}{3.2}
  \clipD{0}{3.2}{0}{3.2}{\flechas{\x}{0}{0.2,0.575,...,3.1}{0.2,0.575,...,3.1}{0.30}{3.1}}
\end{tikzpicture}\\[3pt]{\small $\vc{F}_2=x\,\ux$}
\end{minipage}\hfill
\begin{minipage}[t]{0.325\textwidth}\centering
\begin{tikzpicture}[x=1.12cm, y=1.12cm, baseline=(current bounding box.north)]
  \marcoD{0}{3.2}{0}{3.2}{0,1,2,3}{0,1,2,3}
  \ejesD{0}{3.2}{0}{3.2}
  \clipD{0}{3.2}{0}{3.2}{\flechas{\y}{0}{0.2,0.575,...,3.1}{0.2,0.575,...,3.1}{0.30}{3.1}}
\end{tikzpicture}\\[3pt]{\small $\vc{F}_3=y\,\ux$}
\end{minipage}
\caption{Tres campos vectoriales planos. $\vc{F}_1=\ux+\uy$ es uniforme:
ni diverge ni rota. En $\vc{F}_2=x\,\ux$ el módulo crece \emph{en la dirección
del campo}, lo que da $\nab\cdot\vc{F}_2=1>0$ (expansión) y
$\nab\times\vc{F}_2=\vc{0}$. En $\vc{F}_3=y\,\ux$ el módulo crece
\emph{perpendicularmente} al campo: $\nab\cdot\vc{F}_3=0$ pero
$\nab\times\vc{F}_3=-\uz$ (cizalla). El color codifica $\norm{\vc{F}}$.}
\label{fig:campos_div}
\end{figure}


% -------------------------------------------------------------------------
\subsection{Rotacional}\label{sec:rotacional}
% -------------------------------------------------------------------------
\begin{definicion}[Rotacional]\label{def:rotacional}
El \emph{rotacional} de un campo vectorial derivable $\vc{F}$ es el campo
vectorial
\begin{equation}\label{eq:rot}
  \nab\times\vc{F}\equiv\rot\vc{F}=
  \begin{vmatrix}
    \ux & \uy & \uz\\[2pt]
    \partial_x & \partial_y & \partial_z\\[2pt]
    F_x & F_y & F_z
  \end{vmatrix}
  =\Bigl(\pd{F_z}{y}-\pd{F_y}{z}\Bigr)\ux
  +\Bigl(\pd{F_x}{z}-\pd{F_z}{x}\Bigr)\uy
  +\Bigl(\pd{F_y}{x}-\pd{F_x}{y}\Bigr)\uz ,
\end{equation}
donde el determinante se desarrolla por la primera fila haciendo que las
derivadas (a la izquierda) actúen sobre las componentes (a la derecha). Un campo
de rotacional nulo se llama \emph{irrotacional}.
\end{definicion}

\emph{¿Qué mide?} Si sumergimos en el campo una pequeña ruedecita con
paletas, el rotacional mide la tendencia del campo a hacerla girar:
la \emph{dirección} de $\nab\times\vc{F}$ es el eje de giro (con sentido dado
por la regla de la mano derecha) y su \emph{módulo} es proporcional a la
velocidad angular de la rueda. La definición precisa como densidad de
circulación se da en la \cref{sec:stokes}.

\begin{ejemplo}[Rotacional de campos de referencia]\label{ej:rot}
\begin{enumerate}[label=(\roman*)]
  \item \emph{Giro rígido.} $\vc{F}(\vc{r})=-y\,\ux+x\,\uy$. Solo sobrevive la
        componente $z$ de la \cref{eq:rot}:
        \[
          \nab\times\vc{F}=\Bigl[\pd{(x)}{x}-\pd{(-y)}{y}\Bigr]\uz=[1-(-1)]\,\uz=2\,\uz .
        \]
        El campo «gira» en sentido antihorario, y el rotacional es un vector
        que apunta a lo largo del eje de giro (\cref{fig:campos_rot}, izquierda).
  \item \emph{Cizalla.} $\vc{F}(\vc{r})=y\,\ux$:
        \[
          \nab\times\vc{F}=\Bigl[\pd{(0)}{x}-\pd{(y)}{y}\Bigr]\uz=-\uz .
        \]
        No hay líneas de corriente cerradas, pero el módulo del campo crece en
        la dirección \emph{perpendicular} al campo, y eso basta para hacer girar
        la rueda (en sentido horario).
  \item \emph{Campo radial.} $\nab\times\vc{r}=\vc{0}$: es irrotacional (en
        efecto, $\vc{r}=\nab(r^{2}/2)$ es un gradiente; véase la
        \cref{prop:fundamentales}).
\end{enumerate}
\end{ejemplo}

\begin{observacion}[Cuándo hay rotacional]\label{obs:cuando-rot}
Complementariamente a la \cref{obs:cuando-div}: el rotacional es no nulo
cuando el módulo de una componente del campo cambia en una dirección
\emph{perpendicular} a ella, o cuando el campo gira alrededor de un eje.
\end{observacion}

% Figura: campos con rotacional: F=-y ux + x uy (giro) y F=y ux (cizalla).
\begin{figure}[ht]
\centering
\begin{minipage}[t]{0.47\textwidth}\centering
\begin{tikzpicture}[x=0.58cm, y=0.58cm, baseline=(current bounding box.north)]
  \marcoD{-2.7}{2.7}{-2.7}{2.7}{-2,-1,0,1,2}{-2,-1,0,1,2}
  \ejesD{-2.7}{2.7}{-2.7}{2.7}
  \clipD{-2.7}{2.7}{-2.7}{2.7}{\flechas{-\y}{\x}{-2.4,-1.8,...,2.4}{-2.4,-1.8,...,2.4}{0.34}{3.4}}
\end{tikzpicture}\\[3pt]{\small $\vc{F}(\vc{r})=-y\,\ux+x\,\uy$}
\end{minipage}\hfill
\begin{minipage}[t]{0.47\textwidth}\centering
\begin{tikzpicture}[x=0.58cm, y=0.58cm, baseline=(current bounding box.north)]
  \marcoD{-2.7}{2.7}{-2.7}{2.7}{-2,-1,0,1,2}{-2,-1,0,1,2}
  \ejesD{-2.7}{2.7}{-2.7}{2.7}
  \clipD{-2.7}{2.7}{-2.7}{2.7}{\flechas{\y}{0}{-2.4,-1.8,...,2.4}{-2.4,-1.8,...,2.4}{0.34}{2.5}}
\end{tikzpicture}\\[3pt]{\small $\vc{F}(\vc{r})=y\,\ux$}
\end{minipage}
\caption{Campos con rotacional no nulo. (Izquierda) Giro rígido
$\vc{F}=-y\,\ux+x\,\uy$ (sentido antihorario): $\nab\times\vc{F}=2\,\uz$. Una
pequeña rueda colocada en cualquier punto gira con velocidad angular
constante. (Derecha) Cizalla $\vc{F}=y\,\ux$, de rotacional $-\uz$: no hay
líneas de corriente cerradas, pero la rueda gira igualmente (sentido horario)
porque el flujo es más rápido por un lado que por el otro.}
\label{fig:campos_rot}
\end{figure}


% -------------------------------------------------------------------------
\subsection{Derivadas de productos}\label{sec:productos}
% -------------------------------------------------------------------------
Al tratarse de derivadas, cabe esperar versiones de la regla del producto.
Con dos campos escalares $f,g$ y dos vectoriales $\vc{A},\vc{B}$ se pueden
formar dos productos que dan un escalar ($fg$ y $\vc{A}\cdot\vc{B}$) y dos que
dan un vector ($f\vc{A}$ y $\vc{A}\times\vc{B}$), y sobre cada uno actúa el
operador $\nab$ de la manera que corresponde a su carácter (gradiente sobre
escalares, divergencia y rotacional sobre vectores).

\begin{notacion}[El operador $\vc{A}\cdot\nab$]\label{not:A-nabla}
Dado un campo vectorial $\vc{A}$, el operador escalar
$\vc{A}\cdot\nab=A_x\,\partial_x+A_y\,\partial_y+A_z\,\partial_z$ puede aplicarse
a un campo escalar o a un campo vectorial (componente a componente). Es distinto
de $\nab\cdot\vc{A}$ (\cref{obs:nabla}).
\end{notacion}

\begin{proposicion}[Reglas de derivación del producto]\label{prop:productos}
Sean $f,g$ campos escalares y $\vc{A},\vc{B}$ campos vectoriales de clase
$C^{1}$. Entonces
\begin{align}
  \nab(fg)&=f\,\nab g+g\,\nab f, \label{eq:prod-a}\\
  \nab(\vc{A}\cdot\vc{B})&=\vc{A}\times(\nab\times\vc{B})+\vc{B}\times(\nab\times\vc{A})
     +(\vc{A}\cdot\nab)\vc{B}+(\vc{B}\cdot\nab)\vc{A}, \label{eq:prod-b}\\
  \nab\cdot(f\vc{A})&=f\,\nab\cdot\vc{A}+\vc{A}\cdot\nab f, \label{eq:prod-c}\\
  \nab\cdot(\vc{A}\times\vc{B})&=\vc{B}\cdot(\nab\times\vc{A})-\vc{A}\cdot(\nab\times\vc{B}), \label{eq:prod-d}\\
  \nab\times(f\vc{A})&=\nab f\times\vc{A}+f\,\nab\times\vc{A}, \label{eq:prod-e}\\
  \nab\times(\vc{A}\times\vc{B})&=(\nab\cdot\vc{B})\vc{A}-(\nab\cdot\vc{A})\vc{B}
     +(\vc{B}\cdot\nab)\vc{A}-(\vc{A}\cdot\nab)\vc{B}. \label{eq:prod-f}
\end{align}
\end{proposicion}
\begin{proof}[Demostración de \cref{eq:prod-a,eq:prod-c}]
La primera es la regla de Leibniz aplicada a cada componente: $\partial_i(fg)=f\,\partial_ig+g\,\partial_if$.
Para la segunda, desarrollamos la divergencia en coordenadas y usamos la regla
del producto en cada sumando:
\begin{align*}
  \nab\cdot(f\vc{A})&=\pd{}{x}(fA_x)+\pd{}{y}(fA_y)+\pd{}{z}(fA_z)\\
  &=f\Bigl(\pd{A_x}{x}+\pd{A_y}{y}+\pd{A_z}{z}\Bigr)
   +A_x\pd{f}{x}+A_y\pd{f}{y}+A_z\pd{f}{z}
   =f\,\nab\cdot\vc{A}+\vc{A}\cdot\nab f.
\end{align*}
Las identidades \eqref{eq:prod-d} y \eqref{eq:prod-e} se demuestran con el tensor de
Levi-Civita en el \cref{sec:levi-civita}
(\cref{ejer:div-prod-vect,ejer:rot-fA}); \eqref{eq:prod-b} y \eqref{eq:prod-f} se
obtienen de la misma manera con la identidad $\varepsilon$--$\delta$
(\cref{lema:eps-delta}).
\end{proof}

\begin{observacion}\label{obs:orden}
En las \cref{eq:prod-b,eq:prod-c,eq:prod-f} importa el \emph{orden} de los factores:
$\vc{A}\cdot\nab f$ es la derivada direccional de $f$ a lo largo de $\vc{A}$ (un
escalar), mientras que $\nab\cdot\vc{A}$ es la divergencia. Del mismo modo,
$\vc{B}\cdot\nab$ es un \emph{operador} que se aplica a lo que tiene a su
derecha.
\end{observacion}

% -------------------------------------------------------------------------
\subsection{Operadores de segundo orden}\label{sec:segundo-orden}
% -------------------------------------------------------------------------
Aplicando dos veces $\nab$ se obtienen cinco combinaciones. La más importante
en física es el laplaciano.

\begin{definicion}[Laplaciano]\label{def:laplaciano}
El \emph{laplaciano} de un campo escalar $f$ de clase $C^{2}$ es la divergencia
de su gradiente,
\begin{equation}\label{eq:laplaciano}
  \Delta f\defeq\nab^{2}f=\nab\cdot(\nab f)=\operatorname{div}(\operatorname{grad}f)
  =\sum_{i=1}^{n}\pdd{f}{x_i}\ ;\qquad
  \text{en }\R^{3}:\ \ \Lap f=\pdd{f}{x}+\pdd{f}{y}+\pdd{f}{z}.
\end{equation}
El \emph{laplaciano vectorial} de un campo $\vc{A}$ se define, en coordenadas
cartesianas, componente a componente:
\begin{equation}\label{eq:laplaciano-vect}
  \Lap\vc{A}\defeq(\Lap A_x)\,\ux+(\Lap A_y)\,\uy+(\Lap A_z)\,\uz .
\end{equation}
\end{definicion}

\begin{observacion}
La definición \eqref{eq:laplaciano-vect} es válida \emph{solo} en cartesianas.
En otros sistemas de coordenadas el laplaciano vectorial no es el laplaciano
escalar de cada componente, porque los versores $\ur,\uth,\uph,\ldots$
cambian de punto a punto (\cref{sec:coordenadas}); su definición general es
$\Lap\vc{A}=\nab(\nab\cdot\vc{A})-\nab\times(\nab\times\vc{A})$.
\end{observacion}

Las otras cuatro combinaciones se reducen a dos identidades fundamentales y a
una relación entre ellas.

\begin{proposicion}[Relaciones fundamentales]\label{prop:fundamentales}
Para un campo escalar $f$ y un campo vectorial $\vc{A}$, ambos de clase $C^{2}$,
\begin{equation}\label{eq:rotgrad}
  \nab\times(\nab f)=\vc{0},
  \qquad\qquad
  \nab\cdot(\nab\times\vc{A})=0 .
\end{equation}
\end{proposicion}
\begin{proof}
Ambas se reducen a la simetría de las derivadas parciales segundas
cruzadas (teorema de Schwarz), válida para funciones de clase $C^{2}$.
Desarrollando el rotacional del gradiente,
\[
  \nab\times\nab f=
  \begin{vmatrix}
    \ux & \uy & \uz\\
    \partial_x & \partial_y & \partial_z\\
    \partial_xf & \partial_yf & \partial_zf
  \end{vmatrix}
  =\Bigl(\pdm{f}{y}{z}-\pdm{f}{z}{y}\Bigr)\ux
   +\Bigl(\pdm{f}{z}{x}-\pdm{f}{x}{z}\Bigr)\uy
   +\Bigl(\pdm{f}{x}{y}-\pdm{f}{y}{x}\Bigr)\uz=\vc{0}.
\]
Del mismo modo, para $\vc{A}=(A_x,A_y,A_z)$,
\begin{align*}
  \nab\cdot(\nab\times\vc{A})
  &=\pd{}{x}\Bigl(\pd{A_z}{y}-\pd{A_y}{z}\Bigr)
   +\pd{}{y}\Bigl(\pd{A_x}{z}-\pd{A_z}{x}\Bigr)
   +\pd{}{z}\Bigl(\pd{A_y}{x}-\pd{A_x}{y}\Bigr)\\
  &=\Bigl(\pdm{A_z}{x}{y}-\pdm{A_z}{y}{x}\Bigr)
   +\Bigl(\pdm{A_x}{y}{z}-\pdm{A_x}{z}{y}\Bigr)
   +\Bigl(\pdm{A_y}{z}{x}-\pdm{A_y}{x}{z}\Bigr)=0 .
\end{align*}
En el \cref{sec:levi-civita} se da otra demostración, con índices, que
explica por qué ocurre: un objeto simétrico contraído con otro antisimétrico
es siempre cero.
\end{proof}

\begin{observacion}[Reglas mnemotécnicas, con cautela]\label{obs:mnemo}
El producto vectorial cumple $\vc{v}\times\vc{v}=\vc{0}$, lo que sugiere
$\nab\times\nab=\vc{0}$ y, aplicado a $f$, la primera relación. Cumple también
$\vc{v}\cdot(\vc{v}\times\vc{w})=0$ (el producto vectorial es ortogonal a sus
factores), lo que sugiere la segunda. Pero, como $\nab$ no es un vector, esta
manipulación \emph{sugiere} sin \emph{demostrar}, y puede llevar a error:
aunque $(\vc{v}\times\vc{w})\cdot\vc{w}=0$ para vectores, en general
$(\nab\times\vc{F})\cdot\vc{F}\neq0$.
\end{observacion}

\begin{corolario}[Los gradientes son irrotacionales]\label{cor:grad-irrot}
Si $\vc{F}=\nab f$ es un campo gradiente de clase $C^{1}$ ($f\in C^{2}$), entonces
$\nab\times\vc{F}=\vc{0}$. En particular, ser irrotacional es una condición
\emph{necesaria} para ser conservativo. Es también suficiente en regiones
estrelladas (lema de Poincaré)~\cite{NAVAS}.
\end{corolario}

La combinación restante que se usa en electromagnetismo es la del doble
rotacional.

\begin{proposicion}[Rotacional del rotacional]\label{prop:rotrot}
Para un campo vectorial $\vc{A}$ de clase $C^{2}$,
\begin{equation}\label{eq:rotrot}
  \nab\times(\nab\times\vc{A})=\nab(\nab\cdot\vc{A})-\Lap\vc{A}.
\end{equation}
\end{proposicion}
\begin{proof}
Véase el \cref{ejer:rotrot}, donde se obtiene con el tensor de Levi-Civita y la
identidad $\varepsilon$--$\delta$.
\end{proof}

El resto de combinaciones, $\nab(\nab\cdot\vc{A})$, no tiene un nombre ni un uso
especial en este curso; solo aparece en la \cref{eq:rotrot}.
